With the rise of popularity of chess, areas such as game analysis and engine training, 
have gained increasing academic attention, as a result improved tooling may ease and enhance the tactical and strategic studies of the game.
Chess, with its well-defined structure and semi standardized pieces, presents an ideal 
test bed for non trivial computer vision tasks. \newline
This paper explores a systematic approach to the conversion from images of physical chessboards into the digital FEN representation of the captured game state, 
utilizing a pipeline spanning from a novel chessboard corners detection algorithm to chess-pieces object detection 
using deep learning, ending with the retrieval of similar reached positions. \newline
Our method is designed to handle the challenges posed by the variability of chessboard photos, including different lighting conditions, angles, and chessboard designs.
\newline
We have been through a series of steps to achieve our goal:
\begin{itemize}
    \item We started with the chessboard corners detection, in which we explored different methods both based on deep learning and fully geometric to accurately identify the keypoints on the chessboard.
    \item Next, we focused on chess-pieces object detection, approaching DETR and YOLO to detect and classify chess pieces in various conditions.
    \item Finally, we implemented an image similarity search method using one hot encoding and a similarity search algorithm.
\end{itemize}

The rest of the paper is organized as follows: 
\begin{itemize}
    \item Section I presents our proposed methods for detecting chessboard corners and cell boundaries,.
    \item Section II describes the chess piece detection pipeline.
    \item Section III outlines the procedure for generating FEN annotations from board states.
    \item Section IV introduces the method for chess image similarity search.
\end{itemize}